\section{The Growing Mesh, and How to Picture It}
\label{sec:growth}

The crystal is not static. It grows. New vibes are continually produced, and the mesh expands.
A natural question is where the new vibes go, and the honest answer dissolves a common
confusion.

\paragraph{The mesh is relational, not spatial.}
There is no pre-existing space that the vibes are placed into. In the model, space is the
large-scale shape of the relations, not a container. A vibe has no coordinates. The only sense
in which it is somewhere is what it is connected to. So the question ``where does a new vibe
go?'' has a precise answer: it does not go anywhere. It comes into being attached, by notes, to
some existing vibes, and those attachments are its entire location. A new vibe is a knot tied
onto the existing web, and its position is just which threads it is tied to.

\paragraph{New vibes attach at the present.}
They attach to the frontier, the most recently added vibes, which is the present moment. The
mesh is a causal structure, and it grows the way causal order grows: each new vibe takes some
recent vibes as its past, and is itself part of the past for those that come next. The deep
interior, the distant past, is fixed and never edited. All the activity is at the rim. This
one-way accumulation, relations piling up and never coming undone, is the arrow of time in the
model. The present and the growing edge of the mesh are the same thing.

\paragraph{The self-generation is not a paradox.}
Vibes producing more vibes can sound like something from nothing, but it is ordinary recursion.
A small fixed rule, applied to a seed, unfolds an unlimited structure, in the same way that
``add one'' unfolds all the whole numbers from one. The mesh grows by such a rule, from a seed,
with no new information injected and nothing left to chance.

\paragraph{The canonical image.}
The clearest mental picture is the Poincar\'e disk, the standard way to draw all of hyperbolic
space inside a circle, the device behind Escher's prints of figures shrinking toward a rim. The
center is the origin, the beginning. The radial direction outward is time. Each ring is a
generation of vibes, exponentially more than the ring inside it. The rim is the present, the
growing frontier, and because in hyperbolic geometry the rim is infinitely far away, there is
always room for the next generation. The branching is tree-like and self-similar, so a close
look at any patch shows the same kind of structure again.

Three cautions keep the picture honest. It is a two-dimensional cartoon of a higher-dimensional
object. The real spatial mesh is the three-dimensional honeycomb, and with time it is
four-dimensional. The disk draws relations, not a room, and moving the layout changes nothing,
because the physics is in the connections. And whether one says the infinite crystal already
exists and growth reveals more of it, or that each vibe is freshly created at the frontier, is
a choice of words, since the model commits only to the relations and the rule.

\paragraph{The one-line answer.}
New vibes do not go to a place. They come into being tied by notes to the present frontier of
the web, which, drawn, looks like fresh cells appearing at the ever-receding rim of an
expanding hyperbolic disk, growing a self-similar tree from a single seed.
